\scanpage{7}
9) Le cône sur une variété compacte $S$ peut être une variété à bord, sans que $S$ soit une sphère.

Exemple : considérons le ss-groupe $\Gamma\subset S^3$ d'ordre $120$, \deleted{\unsure{image réciproque}} de $\underset{\substack{\Vert\\ A}}{\unsure{\mathfrak S}_5^+}\subset\underbrace{S^3/\mathord{\pm}1}_{SO(3,\mathbb R)}$. Considérons
\[
S^3/\Gamma\simeq SO(3,\mathbb R)/A\simeq\text{variété des icosaèdres}
\]
\inserted{réguliers} inscrits dans une sphère donnée. Ce n'est pas une sphère, ni sa suspension, qui n'est pas une variété (elle a deux pts singuliers, en lesquels le groupe fondamental local est $\simeq\Gamma$, donc non trivial) Mais \deleted{elle} \inserted{c'est un} \inserted{espace} simplement connexe $S$. \deleted{Sa suspension} \deleted{est encore} Le produit par une droite est néanmoins une variété (?), et sa suspension est une variété \inserted{et $\simeq S^5$,} -- y compris aux deux ``pôles'' ....

\margin{Edwards\par \uline{Exposé Bourbaki}\par 1977 ou 78\par (\unsure{Latour|Lalonde}, 515)}
\margin{$X$ variété de dim $n$, $\mathbb Z$-\unsure{sphère} homologique, sa double susp. est une sphère}

Pour donner des conditions \deleted{topologiques} \inserted{de lissité de nature $\neq$ topolog.} sur un espace à stratification \unsure{str\textsuperscript{t}} conique, \deleted{\unsure{qui impliqueraient que}} \inserted{qui \unsure{viseraient à} \deleted{\illegible{2}}} faire que certains espaces soient des sphères \deleted{topologiques}, il faudrait d'après Deligne se placer dans un cadre

\mspage{6}
